\documentclass[11pt]{article}
\usepackage[a4paper,margin=1in]{geometry}
\usepackage{fontspec}
\usepackage[boldfont=false]{xeCJK}
\setCJKmainfont{FandolSong-Regular.otf}
\usepackage{amsmath}
\usepackage{amssymb}
\usepackage{array}
\usepackage{booktabs}
\usepackage{graphicx}
\usepackage{adjustbox}
\usepackage{enumitem}
\setlist[itemize]{topsep=2pt,partopsep=0pt,parsep=0pt,itemsep=2pt,leftmargin=1.4em}
\setlist[enumerate]{topsep=2pt,partopsep=0pt,parsep=0pt,itemsep=2pt,leftmargin=1.6em}
\usepackage{parskip}
\usepackage{fvextra}
\fvset{breaklines=true,breakanywhere=true,fontsize=\small}
\setlength{\parindent}{0pt}

\begin{document}

\begin{center}
  {\LARGE\bfseries Chemistry of the environment}\\[0.35em]
  {\large Igcse Chemistry}
\end{center}
\vspace{0.6em}

\section*{Water}

\subsection*{Testing for water}

Two chemical tests show that water is \textbf{present}:

\begin{itemize}
\item \textbf{Anhydrous} 无水 cobalt(II) chloride turns from blue to pink when water is added.

\item Anhydrous copper(II) sulfate turns from white to blue when water is added.

\end{itemize}

These tests only show that water is there. To show that water is \textbf{pure}, you test its \textbf{melting point} 熔点 and \textbf{boiling point} 沸点: pure water melts at exactly $0\,^{\circ}\text{C}$ and boils at exactly $100\,^{\circ}\text{C}$. Any dissolved substance changes these values.

This is why \textbf{distilled water} 蒸馏水 is used in chemistry instead of tap water — it has far fewer chemical \textbf{impurities} 杂质.

\subsection*{What is in natural water}

Water from rivers, lakes and the sea is not pure. It may contain dissolved \textbf{oxygen} 氧气, metal compounds, plastics, \textbf{sewage} 污水, harmful \textbf{microbes} 微生物, and \textbf{nitrates} 硝酸盐 and \textbf{phosphates} 磷酸盐 (which come from fertilisers and detergents).

Some of these are helpful:

\begin{itemize}
\item dissolved oxygen lets \textbf{aquatic life} 水生生物 (fish and plants) breathe;

\item some metal compounds give essential \textbf{minerals} 矿物质.

\end{itemize}

Others are harmful:

\begin{itemize}
\item some metal compounds are \textbf{toxic} 有毒 (poisonous);

\item some plastics harm aquatic life;

\item sewage carries microbes that cause disease;

\item nitrates and phosphates cause \textbf{deoxygenation} 缺氧 (loss of oxygen) in the water, which harms aquatic life.

\end{itemize}

\subsection*{Treating drinking water}

To make water safe to drink, the water supply is treated in steps:

\begin{itemize}
\item \textbf{sedimentation} 沉降 and \textbf{filtration} 过滤 remove solid bits;

\item passing it through \textbf{carbon} 碳 removes bad tastes and smells;

\item \textbf{chlorination} 氯消毒 (adding chlorine) kills harmful microbes.

\end{itemize}

\section*{Fertilisers}

\textbf{Fertilisers} 肥料 are added to soil to help plants grow. Ammonium salts and nitrates are common fertilisers.

NPK fertilisers contain the three elements plants need most: \textbf{nitrogen} 氮气, \textbf{phosphorus} 磷 and \textbf{potassium} 钾.

\section*{Air quality and climate}

Clean, dry air is approximately:

\begin{itemize}
\item 78\% nitrogen ($\text{N}_2$)

\item 21\% oxygen ($\text{O}_2$)

\item the rest is a mixture of \textbf{noble gases} 稀有气体 and carbon dioxide ($\text{CO}_2$).

\end{itemize}

\subsection*{Air pollutants and their sources}

\begin{center}
\adjustbox{max width=\linewidth}{%
\begin{tabular}{ll}
\toprule
\textbf{Pollutant} & \textbf{Main source} \\
\midrule
carbon dioxide ($\text{CO}_2$) & \textbf{complete combustion} 完全燃烧 of carbon-containing fuels \\
\textbf{carbon monoxide} 一氧化碳 and \textbf{particulates} 颗粒物 & \textbf{incomplete combustion} 不完全燃烧 of carbon-containing fuels \\
\textbf{methane} 甲烷 & rotting plants and waste gases from animal digestion \\
\textbf{oxides of nitrogen} 氮氧化物 & car engines \\
sulfur dioxide & burning \textbf{fossil fuels} 化石燃料 that contain \textbf{sulfur} 硫 \\
\bottomrule
\end{tabular}}
\end{center}

\subsection*{Effects of these pollutants}

\begin{itemize}
\item \textbf{Carbon dioxide} 二氧化碳 and methane are \textbf{greenhouse gases} 温室气体: more of them causes \textbf{global warming} 全球变暖, which leads to \textbf{climate change} 气候变化.

\item Carbon monoxide is a toxic gas.

\item Particulates increase the risk of \textbf{respiratory} 呼吸 (breathing) problems and \textbf{cancer} 癌症.

\item Oxides of nitrogen cause \textbf{acid rain} 酸雨, \textbf{photochemical smog} 光化学烟雾 and breathing problems.

\item Sulfur dioxide causes acid rain.

\end{itemize}

\subsection*{Reducing these problems}

To slow \textbf{climate change}: plant trees, farm fewer animals, burn fewer fossil fuels, and use more hydrogen and \textbf{renewable energy} 可再生能源 such as wind and solar power.

To reduce \textbf{acid rain}: fit \textbf{catalytic converters} 催化转化器 in cars, use low-sulfur fuels, and use \textbf{flue gas desulfurisation} 烟气脱硫 with \textbf{calcium oxide} 氧化钙 to remove sulfur dioxide from waste gases.

\subsection*{How greenhouse gases warm the Earth}

Greenhouse gases such as carbon dioxide and methane let sunlight through, but they \textbf{absorb} the \textbf{thermal energy} 热能 given off by the warm Earth, and send some of it back down. This reduces the thermal energy lost to space, so the Earth gets warmer.

In a car engine, the high temperature makes nitrogen and oxygen from the air react to form oxides of nitrogen. A \textbf{catalytic converter} removes them, for example:

\[
2\text{CO} + 2\text{NO} \rightarrow 2\text{CO}_2 + \text{N}_2
\]

\subsection*{Photosynthesis}

\textbf{Photosynthesis} 光合作用 is the opposite of combustion — it removes carbon dioxide from the air. Plants use light energy and \textbf{chlorophyll} 叶绿素 to turn carbon dioxide and water into \textbf{glucose} 葡萄糖 and oxygen:

\[
6\text{CO}_2 + 6\text{H}_2\text{O} \rightarrow \text{C}_6\text{H}_{12}\text{O}_6 + 6\text{O}_2
\]


\end{document}
